
An equivalent description of this last homomorphism is as follows: the hypothesis that @@M0@@ is regular implies that these local rings are factorial (by the Auslander--Buchsbaum theorem, EGA IV 21.11.1), and hence that every invertible sheaf @@M1@@ on @@M2@@ extends to an invertible sheaf @@M3@@ on @@M4@@ (EGA IV 21.6.11). One then checks at once that the image under (7.8.1) of @@M5@@ is @@M6@@, where, of course, one puts @@M7@@.

In fact, using only the hypothesis that the local rings of @@M8@@ are factorial and that @@M9@@ is locally integral (without hypotheses of regularity or separatedness), one shows that @@M10@@ does not depend on the chosen extension @@M11@@ of @@M12@@ to an invertible Module on @@M13@@, so that one obtains a homomorphism (7.8.1) directly in this way. To verify the asserted independence, one is reduced to proving that, if @@M14@@, then @@M15@@. Now the hypothesis means that @@M16@@, with @@M17@@ a divisor on @@M18@@ having support @@M19@@ (cf. EGA IV 21.2.11), and one reduces at once to the case where @@M20@@ is \emph{integral}, hence @@M21@@; but then @@M22@@ (@@M23@@ a uniformizer of @@M24@@), so @@M25@@ and a fortiori @@M26@@.

\sgaStableTarget{sga6.e10.sec.7.9}\subsection*{7.9. Specialization and numerical equivalence.}
Suppose that @@M27@@ is proper, and @@M28@@ and @@M29@@ are regular, and consider the two specialization maps (7.4.1) and (7.7.4):
\sgaStableTarget{sga6.e10.eq.7-9-1}@@M30@@
Since the homomorphism
\sgaStableTarget{sga6.e10.eq.7-9-2}@@M31@@
